\documentclass[11pt]{article}
\usepackage[spanish,es-nodecimaldot,es-tabla]{babel}
\usepackage[utf8]{inputenc}
\usepackage[T1]{fontenc}
\usepackage{lmodern}
\usepackage[a4paper,margin=2.5cm]{geometry}
\usepackage{amsmath,amssymb,amsthm}
\usepackage{booktabs,tabularx}
\usepackage[hidelinks]{hyperref}
\usepackage[style=apa,backend=biber]{biblatex}
\addbibresource{referencias.bib}

\newcommand{\TPR}{\mathrm{TPR}}
\newcommand{\FPR}{\mathrm{FPR}}
\newtheorem{hipotesis}{Hipótesis}
\renewcommand{\thehipotesis}{H\arabic{hipotesis}}

\title{Anteproyecto de trabajo final\\[0.3em]
\large Detección de crisis epilépticas en EEG de superficie mediante análisis espectral y redes neuronales: comparación de optimizadores y funciones de activación}
\author{Maestría en Inteligencia Artificial --- Machine Learning\\Universidad de Palermo}
\date{Noviembre de 2026}

\begin{document}
\maketitle

\section{Problema}

La epilepsia afecta a unos 50 millones de personas en el mundo. El diagnóstico y el seguimiento se basan en el electroencefalograma (EEG), cuya revisión manual es lenta: un registro de 24 horas puede llevar horas de lectura de un especialista. Un detector automático de crisis tiene que ser \emph{sensible} (no perder crisis) y, a la vez, \emph{específico} (generar pocas falsas alarmas por hora), dos objetivos en conflicto.

El trabajo aborda la detección de crisis como un problema de clasificación binaria de ventanas de EEG de 4 segundos, con un fuerte desbalance de clases (las crisis ocupan menos del 1~\% del tiempo registrado).

\section{Objetivos}

\paragraph{General.} Construir y evaluar un detector de crisis epilépticas basado en características espectrales del EEG y redes neuronales, y estudiar de forma controlada cómo influyen el optimizador y la función de activación en el entrenamiento y en el rendimiento.

\paragraph{Específicos.}
\begin{enumerate}
  \item Modelar matemáticamente la señal EEG y fundamentar, con la transformada de Fourier, las características espectrales que distinguen una crisis.
  \item Demostrar que la entropía espectral alcanza su máximo global en la distribución uniforme del simplex (multiplicadores de Lagrange) y verificar empíricamente que disminuye durante las crisis.
  \item Comparar seis optimizadores (SGD, SGD con momento de Nesterov, RMSProp, Adam, AdamW y L-BFGS) en un problema convexo (regresión logística) y en uno no convexo (perceptrón multicapa), incluyendo el método del simplex de Nelder-Mead como referencia sin derivadas.
  \item Comparar cinco funciones de activación (sigmoide, tanh, ReLU, Leaky ReLU y GELU) y su interacción con el optimizador y con la profundidad de la red.
  \item Evaluar el detector con curva ROC, AUC, matriz de confusión, sensibilidad y especificidad, y elegir el umbral de decisión como un problema de optimización con restricciones (criterio de Neyman-Pearson).
  \item Controlar el sobreajuste mediante partición por paciente, regularización y \emph{early stopping}.
\end{enumerate}

\section{Hipótesis}

\begin{hipotesis}
En la regresión logística regularizada (problema estrictamente convexo) todos los optimizadores convergen al mismo óptimo global; solo difieren en la velocidad.
\end{hipotesis}
\begin{hipotesis}
En el perceptrón multicapa (problema no convexo) la variabilidad del resultado entre inicializaciones es significativamente mayor que en el caso convexo, como consecuencia de la existencia de múltiples mínimos locales.
\end{hipotesis}
\begin{hipotesis}
Al aumentar la profundidad, las redes con activaciones saturantes (sigmoide y tanh) se degradan más que las de la familia ReLU y GELU, por desvanecimiento del gradiente.
\end{hipotesis}
\begin{hipotesis}
Los optimizadores adaptativos (RMSProp, Adam) son menos sensibles a la elección de la función de activación que SGD.
\end{hipotesis}
\begin{hipotesis}
La entropía espectral del EEG es significativamente menor durante las crisis que fuera de ellas.
\end{hipotesis}

\section{Datos}

Se usará la base pública \emph{CHB-MIT Scalp EEG Database} \parencite{shoeb2009,goldberger2000}, disponible en PhysioNet: registros de EEG de superficie de 23 casos pediátricos (22 sujetos) del Children's Hospital Boston, muestreados a 256~Hz en montaje bipolar del sistema 10-20, con 198 crisis anotadas por especialistas. Se utilizarán los 18 canales comunes a todos los registros, todos los archivos con crisis y una muestra aleatoria de archivos sin crisis (alrededor de 14~GB).

\section{Metodología}

\begin{enumerate}
  \item \textbf{Preprocesamiento:} filtro pasabanda de 0,5 a 45~Hz, filtro \emph{notch} en 60~Hz y segmentación en ventanas de 4~s.
  \item \textbf{Características espectrales} (método de Welch, \cite{welch1967}): potencia absoluta y relativa en las bandas $\delta,\theta,\alpha,\beta,\gamma$, entropía espectral, frecuencia dominante, frecuencia de borde espectral y coherencia entre canales.
  \item \textbf{Partición por sujeto:} entrenamiento, validación y prueba con pacientes disjuntos (los casos \texttt{chb01} y \texttt{chb21}, que son la misma paciente, siempre juntos).
  \item \textbf{Modelos:} regresión logística (convexa) y perceptrón multicapa (no convexo), con entropía cruzada ponderada por clase.
  \item \textbf{Diseño experimental:}
  \begin{itemize}
    \item Ajuste de la tasa de aprendizaje de cada optimizador con igual presupuesto.
    \item Grilla optimizador $\times$ activación $\times$ semilla: $6 \times 5 \times 5 = 150$ entrenamientos.
    \item Profundidad (2, 4 y 8 capas) $\times$ activación, registrando la norma del gradiente por capa y la fracción de unidades inactivas.
    \item Pruebas estadísticas: Friedman, Wilcoxon con corrección de Holm y ANOVA de dos factores \parencite{demsar2006}.
  \end{itemize}
  \item \textbf{Evaluación:} curva ROC, AUC, curva precisión-recall, matriz de confusión, sensibilidad, especificidad, falsas alarmas por hora y sensibilidad por evento (crisis detectadas).
\end{enumerate}

\section{Fundamentación matemática}

\begin{itemize}
  \item \textbf{Transformada de Fourier:} DFT, teorema de Parseval, ventaneo y fuga espectral, estimador de Welch, coherencia espectral \parencite{oppenheim2010,gonzalez2018}.
  \item \textbf{Simplex y Lagrange:} el espectro normalizado $p$ es un punto del simplex $\Delta^{N-1}$. Con el lagrangiano $\mathcal{L}(p,\lambda) = -\sum_i p_i\log p_i + \lambda(\sum_i p_i - 1)$ se demuestra que el máximo de la entropía es $\log N$ en $p_i = 1/N$, y que por la concavidad estricta de $H$ es el óptimo global y único.
  \item \textbf{Optimalidad local y global:} convexidad de la pérdida logística regularizada (Hessiano definido positivo) frente a la no convexidad de las redes; condiciones de primer y segundo orden; métodos de gradiente, momento, adaptativos, cuasi-Newton y Nelder-Mead \parencite{boyd2004,nocedal2006,nelder1965}.
  \item \textbf{Funciones de activación:} teorema de aproximación universal \parencite{cybenko1989,leshno1993}, análisis del gradiente por la regla de la cadena e inicializaciones de Glorot y He \parencite{glorot2010,he2015}.
  \item \textbf{Decisión con restricciones:} $\max_t \TPR(t)$ sujeto a $\FPR(t) \le \alpha$. El multiplicador de Lagrange óptimo es la pendiente de la ROC en el punto de operación, y coincide con el umbral sobre el cociente de verosimilitudes del lema de Neyman-Pearson \parencite{neyman1933,fawcett2006}. Se fijará $\alpha$ equivalente a una falsa alarma por hora.
\end{itemize}

\section{Herramientas}

Python con NumPy, SciPy, scikit-learn, PyTorch y statsmodels. El trabajo completo se entrega como un notebook de Jupyter, ejecutable también en Google Colab, apoyado en módulos \texttt{.py}. El documento se redacta en \LaTeX{} (Overleaf), con una versión equivalente en Word.

\section{Cronograma}

\begin{tabularx}{\linewidth}{@{}lXX@{}}
\toprule
Semana & Desarrollo & Escritura \\
\midrule
1 (12--18/11) & Descarga, preprocesamiento, características, verificación de H5 & Caps.~2 y 3 \\
2 (19--25/11) & Caso convexo, ajuste de tasas, pipeline de evaluación & Caps.~4 y 7 \\
3 (26/11--2/12) & Grilla optimizador $\times$ activación, experimento de profundidad & Caps.~5 y 6 \\
4 (3--9/12) & Pruebas estadísticas, figuras finales & Caps.~8, 9 y 10 \\
10--12/12 & Revisión & Revisión y formato \\
\bottomrule
\end{tabularx}

\section{Índice tentativo}

\begin{enumerate}
  \item Introducción
  \item La señal EEG y su modelo matemático
  \item Análisis de Fourier y características espectrales
  \item Aprendizaje estadístico y sobreajuste
  \item Modelos y funciones de activación
  \item Optimización
  \item Evaluación de clasificadores
  \item Metodología experimental
  \item Resultados
  \item Discusión y conclusiones
\end{enumerate}

\nocite{bishop2024,goodfellow2016,duda2001,buduma2022}
\printbibliography[title={Bibliografía}]

\end{document}
